\documentclass[11pt,a4paper]{article}
\usepackage[margin=2.5cm]{geometry}
\usepackage{amsmath,amsthm}
\usepackage{newtxtext}
\usepackage{newtxmath}
\usepackage{enumitem}
\usepackage{microtype}
\setlength{\parindent}{0pt}
\setlength{\parskip}{4pt}
\allowdisplaybreaks

\newcommand{\E}{\mathbb{E}}
\renewcommand{\P}{\mathbb{P}}
\newcommand{\R}{\mathbb{R}}
\newcommand{\N}{\mathbb{N}}
\newcommand{\Q}{\mathbb{Q}}
\newcommand{\F}{\mathcal{F}}
\newcommand{\G}{\mathcal{G}}
\newcommand{\Hc}{\mathcal{H}}
\newcommand{\1}{\mathbf{1}}
\newcommand{\Leb}{\mathrm{Leb}}
\newcommand{\sgn}{\operatorname{sgn}}
\newcommand{\prob}[1]{\par\medskip\noindent\textbf{#1}\quad}
\newcommand{\sub}[1]{\par\smallskip\noindent\textbf{#1}\ }
\newcommand{\note}[1]{\par\noindent\emph{Remark.} #1}

\begin{document}

\begin{center}
{\large\textbf{Stochastic Processes --- 1060429}}\\[2pt]
Nawar Sgier, 326586047
\end{center}

\noindent Throughout, a Brownian motion (BM) is standard ($B_0=0$) with continuous paths, and
$\F_t$ denotes the filtration in use. Apparent typos are pointed out where they occur.

\section*{Exercise 2}
\subsection*{Conditional expectation}

\prob{1.}
\sub{(a)} Let $Z=\E[X|Y]$. By conditional Jensen, $Z^2\le\E[X^2|Y]$, so $Z\in L^2$. If
$\E[h(Y)^2]=\infty$ the right-hand side is infinite (as $X\in L^2$). Otherwise $V=Z-h(Y)\in
L^2(\sigma(Y))$, and by Problem 2(b) (the product $V(X-Z)$ is integrable by Cauchy--Schwarz)
\[
\E[V(X-Z)]=\E\big[V\,\E[X-Z\,|\,Y]\big]=0 .
\]
Hence $\E(X-h(Y))^2=\E(X-Z)^2+\E V^2+2\E[V(X-Z)]\ge\E(X-Z)^2$.

\sub{(b)} $L^2(\Omega,\sigma(Y),\P)$ is a closed subspace of $L^2(\Omega,\F,\P)$. As shown, $Z=\E[X|Y]$ lies in
it, and $\E[V(X-Z)]=0$ for every $V\in L^2(\sigma(Y))$, i.e.\ $X-Z\perp L^2(\sigma(Y))$. These two
properties characterise the orthogonal projection. (Equivalently, by (a): every $\sigma(Y)$-measurable $V$ is of the form $h(Y)$, so $Z$ is the closest point.)

\prob{2.}
\sub{(a)} ($\Rightarrow$) Let $Z=\E[X|\G]$. For $W=\1_A$, $A\in\G$, the identity is the definition. By linearity
it holds for simple $W$. For $W$ with $|W|\le c$ a.s., take simple $\G$-measurable $W_n\to W$ with $|W_n|\le c$.
Then $|W_nZ|\le c|Z|$ and $|W_nX|\le c|X|$, and dominated convergence gives $\E[WZ]=\E[WX]$.

($\Leftarrow$) Taking $W=\1_A$ gives $\E[Z\1_A]=\E[X\1_A]$ for $A\in\G$. Integrability of $Z$: with
$W_n=\sgn(Z)\1_{\{|Z|\le n\}}$ we get $\E[|Z|\1_{\{|Z|\le n\}}]=\E[W_nX]\le\E|X|$. Letting $n\to\infty$ gives $\E|Z|\le\E|X|$ by
monotone convergence. Hence $Z=\E[X|\G]$.

\sub{(b)} First let $X,Y\ge0$ and $Z=\E[X|\G]\ge0$. For $A\in\G$, $W=\1_A\min(Y,n)$ is bounded and
$\G$-measurable, so $\E[\1_A\min(Y,n)Z]=\E[\1_A\min(Y,n)X]$. By monotone convergence,
$\E[\1_AYZ]=\E[\1_AYX]<\infty$. So $YZ$ is integrable, $\G$-measurable, and equals $\E[XY|\G]$.

In general write $X=X^+-X^-$ and $Y=Y^+-Y^-$. Each $X^\pm Y^\pm\le|XY|$ is integrable, and the
positive case gives $\E[X^\pm Y^\pm|\G]=Y^\pm\E[X^\pm|\G]$. Summing with signs and using linearity,
\[\E[XY|\G]=(Y^+-Y^-)(\E[X^+|\G]-\E[X^-|\G])=Y\,\E[X|\G].\]

\prob{3.}
$\G$ is generated by the partition $I_1=[0,\frac14]$, $I_2=(\frac14,\frac23]$, $I_3=(\frac23,1]$, so
$\G=\{\bigcup_{i\in S}I_i:S\subseteq\{1,2,3\}\}$. Let $Z=\sum_ia_i\1_{I_i}$ with
$a_i=\Leb(I_i)^{-1}\int_{I_i}\omega^2d\omega$. Then $Z$ is $\G$-measurable, and $\E[Z\1_{I_i}]=\E[X\1_{I_i}]$ for each $i$, hence for
every $A\in\G$ by additivity. By uniqueness $\E[X|\G]=Z$ a.s. The coefficients are
\[
a_1=4\cdot\tfrac{1}{3\cdot64}=\tfrac1{48},\qquad
a_2=\tfrac{12}{5}\cdot\tfrac13\Big(\tfrac{8}{27}-\tfrac1{64}\Big)=\tfrac{12}{5}\cdot\tfrac{485}{5184}=\tfrac{97}{432},\qquad
a_3=3\cdot\tfrac13\Big(1-\tfrac{8}{27}\Big)=\tfrac{19}{27}.
\]

\prob{4.}
\note{$Y$ must be allowed to take the value $+\infty$. For instance, if $\G=\{\emptyset,\Omega\}$ then
$\int_\Omega Y\,d\P=\E X=\infty$ forces $Y\equiv\infty$. We prove the statement for $[0,\infty]$-valued $Y$.}

\emph{Existence.} For $M\in\N$ let $X_M=\min(X,M)$ and $Y_M=\E[X_M|\G]$. Since $0\le X_M\le X_{M+1}$, monotonicity gives
$0\le Y_M\le Y_{M+1}$ a.s. Outside a null set the sequence is non-decreasing, and $Y:=\lim_MY_M\in[0,\infty]$ is
$\G$-measurable. For $A\in\G$, monotone convergence gives
\[\int_AY=\lim_M\int_AY_M=\lim_M\int_AX_M=\int_AX.\]

\emph{Uniqueness.} Let $Y,Y'$ be two such variables and $A_n=\{Y'>Y,\ Y\le n\}\in\G$. Then
$\int_{A_n}Y'=\int_{A_n}X=\int_{A_n}Y\le n<\infty$, so $\int_{A_n}(Y'-Y)=0$ with $Y'-Y>0$ on $A_n$. Hence
$\P(A_n)=0$, and taking the union over $n$ gives $\P(Y'>Y,\,Y<\infty)=0$. Since $Y'>Y$ forces $Y<\infty$, we get $\P(Y'>Y)=0$.
By symmetry $\P(Y>Y')=0$.

\prob{5.}
Take $\Omega=\{a,b,c\}$, $\F=2^\Omega$, $\P$ uniform, $X=\1_{\{a\}}$, $\G=\sigma(\{a\})$, $\Hc=\sigma(\{c\})$. Then
$\E[X|\G]=X$ and
\[
\E[\E[X|\G]|\Hc]=\tfrac12\1_{\{a,b\}},\qquad
\E[X|\Hc]=\tfrac12\1_{\{a,b\}},\qquad \E[\E[X|\Hc]|\G]=\tfrac12\1_{\{a\}}+\tfrac14\1_{\{b,c\}}.
\]
At $c$ (which has probability $\frac13$) the first equals $0$ and the second equals $\frac14$.

\subsection*{Martingales}

\prob{1.}
($\Rightarrow$) For $A\in\F_s$, $\E[X_t\1_A]=\E[\E[X_t|\F_s]\1_A]\ge\E[X_s\1_A]$.
($\Leftarrow$) Let $Z=\E[X_t|\F_s]-X_s$, which is $\F_s$-measurable with $\E[Z\1_A]\ge0$ for all $A\in\F_s$.
With $A=\{Z<0\}$ we get $\E[Z\1_{\{Z<0\}}]\ge0$, so $\P(Z<0)=0$, i.e.\ $\E[X_t|\F_s]\ge X_s$.

\prob{2.}
(The index set in $(X_t,\F^X_t)_{t\in X}$ should be $T$.) ($\Leftarrow$) Take $\F_t=\F^X_t$.
($\Rightarrow$) Since $X$ is $(\F_t)$-adapted, $\F^X_t\subseteq\F_t$, and $X$ is $\F^X$-adapted and integrable.
For $s<t$, by the tower property,
\[\E[X_t|\F^X_s]=\E[\E[X_t|\F_s]|\F^X_s]=\E[X_s|\F^X_s]=X_s.\]

\prob{3.}
(Doob's theorem; (a) should read $\E|X_{\tau_2}|<\infty$.) Here $\F_\infty\supseteq\F_n$ for all $n$, and the martingale
property gives $X_n=\E[X_\infty|\F_n]$. A stopping time is $\tau:\Omega\to\N_\infty$ with $\{\tau=n\}\in\F_n$, $n\in\N$, and
$\F_\tau=\{A\in\F_\infty:A\cap\{\tau=n\}\in\F_n\ \forall n\in\N\}$.

\emph{Claim:} for every stopping time $\tau$, $X_\tau=\E[X_\infty|\F_\tau]$.

$X_\tau=\sum_{n\in\N_\infty}X_n\1_{\{\tau=n\}}$ is $\F_\tau$-measurable, since
$\{X_\tau\in B\}\cap\{\tau=n\}=\{X_n\in B\}\cap\{\tau=n\}\in\F_n$. By conditional Jensen, $|X_n|\le\E[|X_\infty||\F_n]$, so
\[\E|X_\tau|\le\sum_n\E[|X_\infty|\1_{\{\tau=n\}}]=\E|X_\infty|.\]
For $A\in\F_\tau$, Fubini (justified by the last bound) gives
\[\E[X_\tau\1_A]=\sum_n\E[X_n\1_{A\cap\{\tau=n\}}]=\sum_n\E[X_\infty\1_{A\cap\{\tau=n\}}]=\E[X_\infty\1_A],\]
because $A\cap\{\tau=n\}\in\F_n$ (trivially for $n=\infty$).

\sub{(a)} $\E|X_{\tau_2}|\le\E|X_\infty|<\infty$.

\sub{(b)} If $A\in\F_{\tau_1}$ then $A\cap\{\tau_2\le n\}=A\cap\{\tau_1\le n\}\cap\{\tau_2\le n\}\in\F_n$, so
$\F_{\tau_1}\subseteq\F_{\tau_2}$. By the claim and the tower property,
\[\E[X_{\tau_2}|\F_{\tau_1}]=\E[\E[X_\infty|\F_{\tau_2}]|\F_{\tau_1}]=\E[X_\infty|\F_{\tau_1}]=X_{\tau_1}.\]

\prob{4.}
$\F_\tau=\{A:A\cap\{\tau\le u\}\in\F_u\ \forall u\}$. From discrete time we use: if $(Z_g,\F_g)_{g\in G}$ is a
supermartingale on a discrete set $G\subset[0,\infty)$ and $\sigma\le\rho\le c$ are $G$-valued stopping
times, then $\E[Z_\rho|\F_\sigma]\le Z_\sigma$, and the stopped process $(Z_{g\wedge\rho})_{g\in G}$ is a
supermartingale.

\emph{Lemma.} Let $Y_n$ be integrable and $\Hc_n$-measurable with $\E[Y_m\1_B]\le\E[Y_n\1_B]$ for all $m\le n$ and
$B\in\Hc_n$, and suppose $L=\sup_n\E Y_n<\infty$. Then $(Y_n)$ is uniformly integrable.

\emph{Proof.} Taking $B=\Omega$ shows $\E Y_n\uparrow L$. From $Y_n\ge\E[Y_1|\Hc_n]\ge-\E[Y_1^-|\Hc_n]$ we get $\E Y_n^-\le\E Y_1^-$, so
$\sup_n\E|Y_n|\le L+2\E|Y_1|=:K$. Given $\varepsilon>0$ pick $m$ with $\E Y_m\ge L-\varepsilon$. For $n\ge m$, since
$\{Y_n\le\lambda\},\{Y_n<-\lambda\}\in\Hc_n$,
\begin{align*}
\E[|Y_n|;|Y_n|>\lambda]&=\E Y_n-\E[Y_n;Y_n\le\lambda]-\E[Y_n;Y_n<-\lambda]\\
&\le\E Y_m+\varepsilon-\E[Y_m;Y_n\le\lambda]-\E[Y_m;Y_n<-\lambda]\le\E[|Y_m|;|Y_n|>\lambda]+\varepsilon .
\end{align*}
As $\P(|Y_n|>\lambda)\le K/\lambda$ and $Y_m$ is integrable, the right side is $\le2\varepsilon$ for $\lambda$ large, uniformly in $n\ge m$.
The finitely many $n<m$ are handled trivially. \hfill$\square$

\sub{(a)} Fix $s<t$ and let $G_n=2^{-n}\N\cup\{s,t\}$; these sets are nested. Let $\tau_n=\inf\{g\in G_n:g>\tau\}$, with $\tau_n=\infty$ if
$\tau=\infty$. Then $\{\tau_n\le g\}=\{\tau<g\}\in\F_g$, so $\tau_n$ is a $G_n$-valued stopping time, and
it is also a stopping time for the grid $G_{n'}$, $n'\ge n$. Moreover $\tau<\tau_n\le\tau+2^{-n}$ and $\tau_n\downarrow\tau$.

For $r\in\{s,t\}$, $X_{r\wedge\tau_n}=\sum_{g\in G_n\cap[0,r]}X_g\1_{\{r\wedge\tau_n=g\}}$ is
$\F_r$-measurable. By right-continuity $X_{r\wedge\tau_n}\to X_{r\wedge\tau}$, so $X_{r\wedge\tau}$ is $\F_r$-measurable.
For $u<r$, $\{X_{r\wedge\tau}\in B\}\cap\{\tau\le u\}=\{X_{u\wedge\tau}\in B\}\cap\{\tau\le u\}\in\F_u$, hence
$X_{r\wedge\tau}$ is $\F_{r\wedge\tau}$-measurable.

Since the stopped discrete process is a supermartingale and $s,t\in G_n$, for $A\in\F_s$,
\[
\E[X_{t\wedge\tau_n}\1_A]\le\E[X_{s\wedge\tau_n}\1_A].\tag{$*$}
\]
Let $Y_n=X_{r\wedge\tau_n}$ and let $\Hc_n$ be the $\sigma$-algebra of $r\wedge\tau_n$ on the grid $G_n$. For $m\le n$, $r\wedge\tau_n\le
r\wedge\tau_m\le r$ are $G_n$-stopping times, so $\E[Y_m|\Hc_n]\le Y_n$. Also $\E Y_n\le\E X_0$. By the Lemma,
$(Y_n)$ is UI. Together with a.s.\ convergence this gives $X_{r\wedge\tau_n}\to X_{r\wedge\tau}$ in $L^1$, so
$X_{r\wedge\tau}$ is integrable. Passing to the limit in $(*)$ gives
$\E[X_{t\wedge\tau}\1_A]\le\E[X_{s\wedge\tau}\1_A]$ for $A\in\F_s$. By Problem 1 (with reversed inequality),
$(X_{t\wedge\tau},\F_t)$ is a supermartingale.

Since $\F_{s\wedge\tau}\subseteq\F_s$ (if $A\in\F_{s\wedge\tau}$ then $A=A\cap\{s\wedge\tau\le s\}\in\F_s$),
$(\F_{t\wedge\tau})$ is a filtration to which $X_{t\wedge\tau}$ is adapted. The same inequality for
$A\in\F_{s\wedge\tau}$ shows that $(X_{t\wedge\tau},\F_{t\wedge\tau})$ is a supermartingale.

\sub{(b)} Applying (a) to $X$ and $-X$, $(X_{t\wedge\tau})$ is a martingale, so $\E X_{t\wedge\tau}=\E X_0$.
On $\{\tau<\infty\}$, $X_{t\wedge\tau}=X_\tau$ for $t\ge\tau$. On $\{\tau=\infty\}$, $X_\tau$ must be given a meaning. We
set $X_\tau:=\lim_tX_{t\wedge\tau}$, which exists a.s.\ by Problem 5(a) because $\sup_t\E X_{t\wedge\tau}^+\le\E Y$. (If
$\tau<\infty$ a.s.\ this is not needed.) Then $|X_\tau|\le Y$, so $X_\tau$ is integrable, and dominated convergence gives
$\E X_\tau=\lim_t\E X_{t\wedge\tau}=\E X_0$.

\prob{5.}
\sub{(a)} Let $D$ be the non-negative dyadic rationals and $F_k=D_k\cap[0,k]$. For $a<b$ let
$U_k$ be the number of upcrossings of $[a,b]$ by $(X_t)_{t\in F_k}$. By Doob's upcrossing inequality for the
discrete submartingale on the finite set $F_k$,
\[(b-a)\E U_k\le\E(X_k-a)^+\le\sup_t\E X_t^++|a|.\]
$U_k$ increases to the number $U_D$ of upcrossings by $(X_t)_{t\in D}$, so by monotone convergence $\E U_D<\infty$. Hence a.s.\
$U_D[a,b]<\infty$ for all rational $a<b$, and $\lim_{D\ni t\to\infty}X_t$ exists in $[-\infty,\infty]$ (otherwise
$\liminf<a<b<\limsup$ would give infinitely many upcrossings).

By right-continuity, if $X_t>c$ for arbitrarily large real $t$ then $X_d>c$ for arbitrarily large $d\in D$. So
$\limsup_{t\to\infty}X_t=\limsup_{D\ni t\to\infty}X_t$, and similarly for $\liminf$. Thus $X=\lim_{t\to\infty}X_t$ exists a.s.
Since $\E|X_t|=2\E X_t^+-\E X_t\le2\sup_t\E X_t^+-\E X_0$, Fatou's lemma gives $\E|X|<\infty$.

\sub{(b)} \note{Right-continuity must be assumed here as well. Let $U\sim U[0,1]$ and
$X_t=\1_{\{t-\lfloor t\rfloor=U\}}$. Then $X_t=0$ a.s.\ for each $t$, so $X$ is an $L^2$-bounded martingale, but
every path equals $1$ at $U+k$ for all $k\in\N$ and does not converge.}

Assume right-continuity. Since $\sup\E X_t^+\le\sup(\E X_t^2)^{1/2}<\infty$, (a) gives $X_t\to X$ a.s.
For $s<t$, by Conditional expectation, Problem 2(b), $\E[X_s(X_t-X_s)]=\E[X_s\E[X_t-X_s|\F_s]]=0$. Hence
$\E(X_t-X_s)^2=\E X_t^2-\E X_s^2$. So $\E X_t^2$ is non-decreasing and bounded, and $\E(X_t-X_s)^2\to0$ as $s,t\to\infty$.
By completeness of $L^2$, $X_t\to X'$ in $L^2$. Along a sequence $t_n\to\infty$ a subsequence converges a.s.\ to $X'$,
so $X'=X$ a.s.\ and $X\in L^2$. Finally, for $u>t$, $X_t=\E[X_u|\F_t]$ and
\[\|\E[X_u|\F_t]-\E[X|\F_t]\|_1\le\|X_u-X\|_1\to0,\]
so $\E[X|\F_t]=X_t$ a.s.

\prob{6.}
\sub{(a)} The argument works for every $\sigma\in\R$, which we need in (b). Let $M_t(\sigma)=\exp(\sigma
B_t-\sigma^2t/2)$. It is adapted, and integrable since $\E e^{\sigma B_t}=e^{\sigma^2t/2}$. Since $B_t-B_s$ is independent
of $\F_s$,
\[\E[M_t|\F_s]=M_s\,\E\big[e^{\sigma(B_t-B_s)}\big]e^{-\sigma^2(t-s)/2}=M_s.\]

\sub{(b)} Writing $x=B_t$,
$e^{\sigma x-\sigma^2t/2}=e^{\sigma x}e^{-\sigma^2t/2}=\sum_k\frac{\sigma^k}{k!}H_k(x,t)$, where
\[
H_2=x^2-t,\quad H_3=x^3-3tx,\quad H_4=x^4-6tx^2+3t^2
\]
are read off from the coefficients of $\sigma^2,\sigma^3,\sigma^4$ in
$(1+\sigma x+\frac{\sigma^2x^2}2+\frac{\sigma^3x^3}6+\frac{\sigma^4x^4}{24})(1-\frac{\sigma^2t}2+\frac{\sigma^4t^2}8)$.
Thus $\partial_\sigma^kM_t(\sigma)|_{\sigma=0}=H_k(B_t,t)$.

By (a), $\E[M_t(\sigma)\1_A]=\E[M_s(\sigma)\1_A]$ for $A\in\F_s$ and all $\sigma$. For $|\sigma|\le1$,
$\partial^j_\sigma M_t(\sigma)=P_j(B_t-\sigma t,t)M_t(\sigma)$ with $P_j$ a polynomial. This is dominated by
$C_j(1+|B_t|+t)^je^{|B_t|}$, which is integrable, so we may differentiate under $\E$ $k$ times at $\sigma=0$:
$\E[H_k(B_t,t)\1_A]=\E[H_k(B_s,s)\1_A]$. Since $H_k(B_t,t)$ is adapted and integrable, it is a martingale. For
$k=2,3,4$ this gives the processes $B_t^2-t$, $B_t^3-3tB_t$ and $B_t^4-6tB_t^2+3t^2$.


\end{document}
